\themaN
\graphicspath{{../Ch24_Comparaison_de_fractions/Images/}}

\chapter{Calcul avec\\des fractions}
\label{C20}


%%%%%%%%%%%%%%%%%%%%%%%%%%%%%%%%%
%%%%%%%%%%%%%%%%%%%%%%%%%%%%%%%%%
\begin{prerequis}[Connaissances et compétences abordées]
   \begin{itemize}
      \item Encadrer une fraction par deux nombres entiers consécutifs.
      \item Comparer deux fractions de même dénominateur.
      \item Ecrire une fraction comme somme d’un entier et d’une fraction inférieure à 1.
      \item Connaître des égalités entre des fractions usuelles (exemples : 5/10 = 1/2 ; 10/100 = 1/10 ; 2/4 = 1/2).
   \end{itemize}
\end{prerequis}

\vfill

\begin{debat}[Débat : l'\oe il d'Oudjat, ou \oe il d'Horus] 
   Dans la mythologie égyptienne, on trouve une histoire liée aux fractions : {\it Seth}, dieu de la violence et incarnation du mal, arrache l’{\bf œil} à {\bf Horus}. {\it Seth} le partage en six morceaux et les répand à travers l’Egypte. {\it Thot}, dieu magicien, reconstitue l’œil, symbole du bien contre le mal. Chacune de ses parties symbolise une fraction de numérateur 1 et de dénominateurs 2, 4, 8, 16, 32 et 64. Il accordera le 64\up{e} manquant à tout scribe recherchant et acceptant sa protection.
   \begin{center}
      \includegraphics[width=6cm]{Horus}
   \end{center}
   \bigskip
   \begin{cadre}[B2][F4]
      \begin{center}
         Vidéo : \href{https://www.youtube.com/watch?v=yHft4m1Gi7k}{\bf La légende de l'\oe il d'Horus}, chaîne YouTube {\it Ma deuxième école}, épisode de la série {\it Culture G}.
      \end{center}
   \end{cadre}
\end{debat}

\vfill

\textcolor{PartieGeometrie}{\sffamily\bfseries Cahier de compétences} : chapitre 4, exercices 16 à 30 ; 37 à 42 ; 48 à 50.


%%%%%%%%%%%%%%%%%%%%%%%%%%%%%%%%%%%%
%%%%%%%%%%%%%%%%%%%%%%%%%%%%%%%%%%%%
\activites

\begin{activite}[Décomposer une fraction]
   {\bf Objectifs :} représenter des fractions ; écrire une fraction sous la forme de la somme d'un entier et d'une fraction inférieure à 1.
   \begin{QCM}
      \partie[des toasts en entrée]
         Pour faire des toasts à ses amis, Mattis coupe des tranches de pain de mie en quatre, avant de les garnir. \\
         Samy mange onze de ces petits toasts. Quelle fraction de grande tranche de pain de mie Samy a-t-il mangée ?
         \begin{center}
            {\psset{xunit=0.9,yunit=0.8}
            \begin{pspicture}(0,0)(3,1.5)
               \psframe[fillstyle=solid,fillcolor=B3](0,0)(2,1.5)
               \psline(1,0)(1,1.5)
               \psline(0,0.75)(2,0.75)
            \end{pspicture}
            \begin{pspicture}[fillstyle=solid,fillcolor=B3](0,0)(3,1.5)
               \psframe(0,0)(2,1.5)
               \psline(1,0)(1,1.5)
               \psline(0,0.75)(2,0.75)
            \end{pspicture}
            \begin{pspicture}(0,0)(3,1.5)
               \psframe(0,0)(2,1.5)
               \pspolygon[fillstyle=solid,fillcolor=B3](0,0)(1,0)(1,0.75)(2,0.75)(2,1.5)(0,1.5)
               \psline(1,0)(1,1.5)
               \psline(0,0.75)(2,0.75)
            \end{pspicture}
            \begin{pspicture}(0,0)(3,1.5)
               \psframe(0,0)(2,1.5)
               \psline(1,0)(1,1.5)
               \psline(0,0.75)(2,0.75)
            \end{pspicture}
            \begin{pspicture}(0,0)(3,1.5)
               \psframe(0,0)(2,1.5)
               \psline(1,0)(1,1.5)
               \psline(0,0.75)(2,0.75)
            \end{pspicture}}
         \end{center}
         Il a mangé 11 fois $\dfrac14$ de grande tranche, donc $\dfrac{11}4$ de grande tranche, ce qui fait 2 tranches et $\dfrac34$ de tranche. \\ [2mm]
On peut écrire : \hspace{40mm} $\dfrac{11}4 =2+\dfrac34$. \\ [2mm]
         Meryem, elle, a mangé 17 petits toasts. Colorier ce que cela représente ci-dessous.
         \begin{center}
            {\psset{xunit=0.9,yunit=0.8}
            \begin{pspicture}(0,0)(3,1.5)
               \psframe(0,0)(2,1.5)
               \psline(1,0)(1,1.5)
               \psline(0,0.75)(2,0.75)
            \end{pspicture}
            \begin{pspicture}(0,0)(3,1.5)
               \psframe(0,0)(2,1.5)
               \psline(1,0)(1,1.5)
               \psline(0,0.75)(2,0.75)
            \end{pspicture}
            \begin{pspicture}(0,0)(3,1.5)
               \psframe(0,0)(2,1.5)
               \pspolygon(0,0)(1,0)(1,0.75)(2,0.75)(2,1.5)(0,1.5)
               \psline(1,0)(1,1.5)
               \psline(0,0.75)(2,0.75)
            \end{pspicture}
            \begin{pspicture}(0,0)(3,1.5)
               \psframe(0,0)(2,1.5)
               \psline(1,0)(1,1.5)
               \psline(0,0.75)(2,0.75)
            \end{pspicture}
            \begin{pspicture}(0,0)(3,1.5)
               \psframe(0,0)(2,1.5)
               \psline(1,0)(1,1.5)
               \psline(0,0.75)(2,0.75)
            \end{pspicture}}
         \end{center} \medskip
         Compléter l'égalité  : \hspace{30mm} $\dfrac{17}{4} = \hdashrule{1.1cm}{0.15pt}{4pt 2pt} + \dfrac{\hdashrule{1cm}{0.15pt}{4pt 2pt}}{\hdashrule{1cm}{0.15pt}{4pt 2pt}}$ \\

      \partie[des \og flam \fg{} en plat principal]
         Pour poursuivre, Mattis propose à ses amis de manger des flammekueches (flam). \\
         \begin{minipage}{7.75cm}
            \begin{pspicture}(-1.2,-1)(1,1)
               \pscircle[fillstyle=solid,fillcolor=J2](0,0){0.7}
               \multido{\n=0+72}{5}{\psline(0,0)(0.7;\n)}
            \end{pspicture}
            \begin{pspicture}(-0.7,-1)(1,1)
               \pscircle(0,0){0.7}
               \pswedge[fillstyle=solid,fillcolor=J2](0,0){0.7}{72}{288}
               \multido{\n=0+72}{5}{\psline(0,0)(0.7;\n)}
            \end{pspicture}
            \begin{pspicture}(-0.7,-1)(1,1)
               \pscircle(0,0){0.7}
               \multido{\n=0+72}{5}{\psline(0,0)(0.7;\n)}
            \end{pspicture}
            \begin{pspicture}(-0.7,-1)(1,1)
               \pscircle(0,0){0.7}
               \multido{\n=0+72}{5}{\psline(0,0)(0.7;\n)}
            \end{pspicture} \\
            Une part de flam représente $\dfrac{\hdashrule{1cm}{0.15pt}{4pt 2pt}}{\hdashrule{1cm}{0.15pt}{4pt 2pt}}$ de flam. \\ [3mm]
            $\dfrac{\hdashrule{1cm}{0.15pt}{4pt 2pt}}{\hdashrule{1cm}{0.15pt}{4pt 2pt}}$ de flam est colorié. \\ [3mm]
            Donc, $\dfrac{\hdashrule{1cm}{0.15pt}{4pt 2pt}}{\hdashrule{1cm}{0.15pt}{4pt 2pt}} ={\hdashrule{1cm}{0.15pt}{4pt 2pt}} + \dfrac{\hdashrule{1cm}{0.15pt}{4pt 2pt}}{\hdashrule{1cm}{0.15pt}{4pt 2pt}}$ \\ [2mm]
         \end{minipage}
         \qquad
         \begin{minipage}{7.75cm}
            \begin{pspicture}(-1.2,-1)(1,1)
               \pscircle(0,0){0.7}
               \multido{\n=0+45}{8}{\psline(0,0)(0.7;\n)}
               \end{pspicture}
               \begin{pspicture}(-0.7,-1)(1,0.8)
               \pscircle(0,0){0.7}
               \multido{\n=0+45}{8}{\psline(0,0)(0.7;\n)}
               \end{pspicture}
               \begin{pspicture}(-0.7,-1)(1,1)
               \pscircle(0,0){0.7}
               \multido{\n=0+45}{8}{\psline(0,0)(0.7;\n)}
               \end{pspicture}
               \begin{pspicture}(-0.7,-1)(1,1)
               \pscircle(0,0){0.7}
               \multido{\n=0+45}{8}{\psline(0,0)(0.7;\n)}
               \end{pspicture} \\
               Une part de flam représente $\dfrac{\hdashrule{1cm}{0.15pt}{4pt 2pt}}{\hdashrule{1cm}{0.15pt}{4pt 2pt}}$ de flam. \\
               Colorier $\dfrac{13}{8}$ de flam. \\ [1mm]
               Donc, $\dfrac{13}{8} ={\hdashrule{1cm}{0.15pt}{4pt 2pt}}+\dfrac{\hdashrule{1cm}{0.15pt}{4pt 2pt}}{\hdashrule{1cm}{0.15pt}{4pt 2pt}}$ \\ [2mm]
         \end{minipage}

      \partie[Des éclairs en dessert]
         Julie a mangé $\dfrac{73}9$ de mini-éclairs au chocolat, trouver une manière de décomposer cette fraction en la somme \\ [1mm]
         d'un nombre entier et d'une fraction inférieure à 1. \\ [22mm]
   \end{QCM}
   \vfill\hfill{\it\footnotesize{Inspiré de \href{http://ww2.ac-poitiers.fr/dsden86-pedagogie/sites/dsden86-pedagogie/IMG/pdf/groupe4_c13.pdf}{Écrire une fraction sous la forme d'un entier et d'une fraction inférieure à 1}, académie de Poitiers.}}
\end{activite}


%%%%%%%%%%%%%%%%%%%%%%%%%%%%%%%%%%
%%%%%%%%%%%%%%%%%%%%%%%%%%%%%%%%%%
\cours 

%%%%%%%%%%%%%%%%%%%%%%%%%%%%%%%%%%
\section{Comparer des fractions}

\begin{methode}[Comparer des fractions de même dénominateur]
   Pour comparer deux fractions ayant le même dénominateur, il suffit de comparer les numérateurs : la fraction ayant le plus grand numérateur est la plus grande.
   \exercice \smallskip
      Comparer $\dfrac26$ et $\dfrac36$.
   \correction \smallskip
      $2<3$ donc $\dfrac26<\dfrac36$
      {\psset{unit=0.6}
      \begin{pspicture}(-1.8,-0.5)(1,0.2)
         \pscircle(0,0){1}
         \pswedge[fillstyle=solid,fillcolor=B2](0,0){1}{60}{180}
         \multido{\n=0+60}{6}{\psline(0,0)(1;\n)}
         \rput(1.25,0){$<$}
      \end{pspicture}
      \begin{pspicture}(-1.5,-0.5)(1,0.2)
         \pscircle(0,0){1}
         \pswedge[fillstyle=solid,fillcolor=B2](0,0){1}{0}{180}
         \multido{\n=0+60}{6}{\psline(0,0)(1;\n)}
      \end{pspicture}}
\end{methode}

\begin{remarque}
   cela revient à dire que lorsque l'on partage une quantité en un nombre de parts égales, plus on prend de parts plus le nombre est grand.
\end{remarque}

\begin{methode}[Comparer une fraction à 1]
   Pour savoir si une fraction $\dfrac{a}{b}$ est inférieure, égale ou supérieure à 1, on compare le numérateur $a$ au dénominateur $b$ :
   \begin{center}
      \begin{tabular}{|c|C{1.5}|C{1.5}|C{1.5}|}
         \hline
         si & $a<b$ & $a=b$ & $a>b$ \\
         \hline
         alors & $\dfrac{a}{b}<1$ & $\dfrac{a}{b}=1$ & $\dfrac{a}{b}>1$ \\ [1mm]
         \hline
      \end{tabular}
   \end{center}
   \exercice \smallskip
      Comparer $\dfrac36$, $\dfrac86$ et $\dfrac66$ à 1.
   \correction
      dans $\dfrac36$, on a $3<6$ donc, $\dfrac36<1$
      {\psset{unit=0.6}
      \begin{pspicture}(-3,-0.5)(1,0.2)
         \pscircle(0,0){1}
         \pswedge[fillstyle=solid,fillcolor=B2](0,0){1}{0}{180}
         \multido{\n=0+60}{6}{\psline(0,0)(0.98;\n)}
      \end{pspicture}} \\
      dans $\dfrac86$, on a $8>6$ donc, $\dfrac86>1$
      {\psset{unit=0.6}
      \begin{pspicture}(-1.5,-0.3)(1,0.8)
         \pscircle[fillstyle=solid,fillcolor=B2](0,0){1}
         \multido{\n=0+60}{6}{\psline(0,0)(0.98;\n)}
      \end{pspicture}
      \begin{pspicture}(-1.9,-0.3)(1,0.8)
         \pscircle(0,0){1}
         \pswedge[fillstyle=solid,fillcolor=B2](0,0){1}{60}{180}
         \multido{\n=0+60}{6}{\psline(0,0)(0.98;\n)}
      \end{pspicture}} \\
      dans $\dfrac66$, on a $6=6$ donc, $\dfrac66=1$
      {\psset{unit=0.6}
      \begin{pspicture}(-3,0)(1,1)
         \pscircle[fillstyle=solid,fillcolor=B2](0,0){1}
         \multido{\n=0+60}{6}{\psline(0,0)(0.98;\n)}
      \end{pspicture}} \\
\end{methode}

%\begin{remarque}
%   cela revient à dire que lorsque l'on partage une quantité en un nombre de parts égales, si on prend un nombre de parts plus petit que le nombre total, on a une quantité plus petite qu'un entier ; si on en prend autant on a la quantité en totalité et si on en prend plus, on a plus qu'un entier.
%\end{remarque}
   

%%%%%%%%%%%%%%%%%%%%%%%%%%%%%
\section{Décomposer une fraction}

\begin{methode}[Décomposer une fraction]
   Pour décomposer une fraction en somme d'un entier et d'une fraction inférieure à 1, on effectue la division euclidienne du numérateur par le dénominateur. Le quotient nous donne le nombre entier et la fraction inférieure à 1 s'obtient en prenant comme numérateur le reste et comme dénominateur le diviseur.
   \exercice \smallskip
   Décomposer $\dfrac{23}{6}$.
   \correction \smallskip
   $\opidiv{23}{6}$ \qquad donc, \qquad $23 =6\times3+5$ \quad et \quad $\dfrac{23}{6} =3+\dfrac56$
\end{methode}

\bigskip

\begin{minipage}{8cm}
   Cette décomposition permet d'encadrer facilement \\
   un nombre entre deux entiers consécutifs : \\ [2mm]
   par exemple, $\dfrac{23}{6} =3+\dfrac56$ donc $3<\dfrac{23}{6}<4$.
\end{minipage}
\begin{minipage}{8cm}
  {\psset{unit=0.6}
  \begin{pspicture}(-2,-1)(1.5,1)
      \pscircle[fillstyle=solid,fillcolor=B2](0,0){1}
      \multido{\n=0+60}{6}{\psline(0,0)(1;\n)}
   \end{pspicture}
   \begin{pspicture}(-1,-1)(1.5,1)
      \pscircle[fillstyle=solid,fillcolor=B2](0,0){1}
      \multido{\n=0+60}{6}{\psline(0,0)(1;\n)}
   \end{pspicture}
   \begin{pspicture}(-1,-1)(1.5,1)
      \pscircle[fillstyle=solid,fillcolor=B2](0,0){1}
      \multido{\n=0+60}{6}{\psline(0,0)(1;\n)}
   \end{pspicture}
   \begin{pspicture}(-1,-1)(1,1)
      \pscircle(0,0){1}
      \pswedge[fillstyle=solid,fillcolor=B2](0,0){1}{60}{360}
      \multido{\n=0+60}{6}{\psline(0,0)(1;\n)}
   \end{pspicture}}
\end{minipage}

  
%%%%%%%%%%%%%%%%%%%%%%%%%%%%%%%%%%
%%%%%%%%%%%%%%%%%%%%%%%%%%%%%%%%%%
\exercicesbase

\begin{colonne*exercice}

%%%%%%%%%%%%%
\serie{Comparer des fractions}

\begin{exercice}
   Comparer les fractions suivantes : \medskip
   \begin{colenumerate}{2}
      \item $\dfrac{1}{9}$ et $\dfrac{12}{9}$ \bigskip
      \item $\dfrac{5}{3}$ et $\dfrac{2}{3}$ \bigskip
      \item $\dfrac{18}{17}$ et $\dfrac{17}{17}$ \bigskip
      \item $\dfrac{7}{129}$ et $\dfrac{17}{129}$
      \item $\dfrac{123}{6}$ et $\dfrac{211}{6}$
      \item $\dfrac{181}{13}$ et $\dfrac{118}{13}$
   \end{colenumerate}
\end{exercice}

\medskip

\begin{exercice}
   Pour chaque fraction, dire si elle est inférieure, égale ou supérieure à 1. \medskip
   \begin{colenumerate}{3}
      \item $\dfrac72$ \bigskip
      \item $\dfrac22$ \bigskip
      \item $\dfrac{11}3$ \bigskip
      \item $\dfrac53$
      \item $\dfrac{122}{132}$
      \item $\dfrac{15}{14}$
      \item $\dfrac{13}{14}$
      \item $\dfrac{14}{14}$
      \item $\dfrac{154}{145}$
   \end{colenumerate}
\end{exercice}

\medskip

\begin{exercice}
   Les élèves font le cross du collège par équipes de classe. Les professeurs d'EPS ont écrit le temps de chaque équipe sous la forme d'une fraction d'heure. \ \medskip
   \begin{colitemize}{3}
      \item 6A : $\dfrac{13}{6}$ h \bigskip
      \item 6B : $\dfrac86$ h \bigskip
      \item 6C : $\dfrac{10}{15}$ h \bigskip
      \item 6D : $\dfrac{45}{60}$ h \bigskip
      \item 6E : $\dfrac{11}3$ h
      \item 6F : $\dfrac73$ h
   \end{colitemize}
   \vspace*{-5mm}
   \begin{enumerate}
      \item Quelles classes ont mis moins d'une heure ?
      \item Quelles classes ont mis entre une et deux heures ?
      \item Quelles classes ont mis plus de deux heures ?
      \item Établir un classement.
   \end{enumerate}
\end{exercice}

\medskip

%%%%%%%%%%%%%
\serie{Décomposer des fractions}

\begin{exercice}
   Écrire chaque fraction sous la forme de la somme d'un nombre entier et d'une fraction inférieure à 1 en faisant le calcul mentalement. \medskip
   \begin{colenumerate}{3}
      \item $\dfrac72$ \bigskip
      \item $\dfrac98$
      \item $\dfrac13$
      \item $\dfrac53$
      \item $\dfrac{11}{4}$
      \item $\dfrac{13}{4}$
   \end{colenumerate}
\end{exercice}

\medskip

\begin{exercice}
   Écrire chaque fraction sous la forme de la somme d'un nombre entier et d'une fraction inférieure à 1 en posant la division euclidienne. \medskip
   \begin{colenumerate}{3}
      \item $\dfrac{27}2$ \bigskip
      \item $\dfrac{92}8$ \bigskip
      \item $\dfrac{13}{3}$
      \item $\dfrac{53}{3}$
      \item $\dfrac{119}{4}$
      \item $\dfrac{127}{4}$
   \end{colenumerate}
\end{exercice}

\medskip

\begin{exercice}
   Encadrer chacune des fractions suivantes par deux nombres entiers consécutifs. \medskip
    \begin{colenumerate}{3}
      \item $\dfrac72$ \bigskip
      \item $\dfrac98$ \bigskip
      \item $\dfrac13$ \bigskip
      \item $\dfrac53$ \bigskip
      \item $\dfrac{11}{4}$
      \item $\dfrac{13}{4}$
      \item $\dfrac{27}2$
      \item $\dfrac{92}8$
      \item $\dfrac{13}{3}$
      \item \, $\dfrac{53}{3}$
      \item \, $\dfrac{119}{4}$
      \item \, $\dfrac{127}{4}$
   \end{colenumerate}
\end{exercice}

\medskip

\begin{exercice}
   Répondre aux petits défis suivante : \medskip
   \begin{enumerate}
      \item Combien faut-il ajouter à $\dfrac23$ pour obtenir 1 ? \bigskip
      \item Combien faut-il ajouter à $\dfrac74$ pour obtenir 2 ? \bigskip
      \item Combien faut-il ajouter à $\dfrac3{10}$ pour obtenir 2 ? \bigskip
      \item Combien faut-il ajouter à $\dfrac{12}7$ pour obtenir 3 ? \bigskip
   \end{enumerate}
\end{exercice}

\medskip

\begin{exercice}
   Relier chaque fraction à une fraction égale. \\ [2mm]
   \begin{tabular}{rcp{2cm}cp{2cm}}
      $\dfrac{5}{10}$ & $\bullet$ & & $\bullet$ & $\dfrac{1}{10}$ \\ [4mm]
      $\dfrac{10}{100}$ & $\bullet$ & & & \\ [4mm]
      $\dfrac{2}{4}$ & $\bullet$ & & $\bullet$ & $\dfrac{1}{5}$ \\ [4mm]
      $\dfrac{100}{1\,000}$ & $\bullet$ & & & \\ [4mm]
      $\dfrac{2}{10}$ & $\bullet$ & & $\bullet$ & $\dfrac{1}{2}$ \\
   \end{tabular}
\end{exercice} 

\end{colonne*exercice}


%%%%%%%%%%%%%%%%%%%%%%%%%%%%%%
%%%%%%%%%%%%%%%%%%%%%%%%%%%%%%
\Recreation

\enigme[Sudofractions]
   \begin{minipage}{7.5cm}
      \begin{itemize}
         \item Pour chaque fraction ou nombre décimal, déterminer le nombre entier immédiatement inférieur.
         \item Placer ce nombre entier dans la grille vierge ci-contre, exactement au même endroit que sur la grille de calculs.
         \item Compléter la grille selon les règles du Sudoku.
      \end{itemize}
      {\it Par exemple, $\dfrac{5}{2} =2+\dfrac{1}{2} (=2,5)$ donc, \\ [1mm]
         la troisième case de la ligne du haut comporte un 2}.
   \end{minipage}
   \qquad
   \begin{minipage}{7cm}
  {\hautab{1.8}
      \begin{tabular}{|*{3}{C{0.5}|}|*{3}{C{0.5}|}|*{3}{C{0.5}|}}
         \hline
         & & {\bf 2} & & & & & & \\
         \hline
         & & & & & & & & \\
         \hline
         & & & & & & & & \\
         \hline
         \hline
         & & & & & & & & \\
         \hline
         & & & & & & & & \\
         \hline
         & & & & & & & & \\
         \hline
         \hline
         & & & & & & & & \\
         \hline
         & & & & & & & & \\
         \hline
         & & & & & & & & \\
         \hline
      \end{tabular}}
   \end{minipage}
  \vfill
   \begin{center}
   {\hautab{3.44}
   \footnotesize
      \begin{tabular}{|*{3}{C{0.9}|}|*{3}{C{0.9}|}|*{3}{C{0.9}|}}
         \hline
          \!$\dfrac{13}{2}$\! & $\dfrac{27}{5}$ & $\dfrac{5}{2}$ & & 8,73 & 7,99 & $\dfrac{39}{4}$ & & \!$\dfrac{43}{10}\!$ \\
         \hline
          $\dfrac{77}{10}$ & $\dfrac{7}{6}$ & & 4,5 & & & $\dfrac{25}{7}$ & & \\
         \hline
          & & $\dfrac{17}{2}$ & $1,001$ & & & & & $\dfrac{38}{5}$\\
          \hline
          \hline
          8,25 & 2,89 & & & $\dfrac{33}{10}$ & & & & $6,7$ \\
          \hline
          5,55 & & & $\dfrac{21}{9}$ & & & & 3,199 & \\
         \hline
          & & & $\dfrac{80}{9}$ & & & 5,356 & 2,1 & \\
         \hline
         \hline
          & $\dfrac{845}{100}$ & $\dfrac{60}{8}$ & 5,2 & 2,8 & & $\dfrac{13}{2}$ & $\dfrac{25}{6}$ & \\
         \hline
          \!$\dfrac{30}{11}$\! & $\dfrac{25}{8}$ & & 6,001 & & 8,81 & & $\dfrac{65}{7}$ & \\
         \hline
          $\dfrac{29}{3}$ & & & & 4,909 & & & & \\
         \hline
      \end{tabular}}
   \end{center}

